\documentclass[11pt]{article}

\usepackage[margin=1in]{geometry}
\usepackage[T1]{fontenc}
\usepackage{lmodern}
\usepackage{microtype}
\usepackage{hyperref}
\usepackage{xcolor}
\usepackage{listings}
\usepackage{enumitem}

\hypersetup{
  colorlinks=true,
  linkcolor=blue!50!black,
  citecolor=blue!50!black,
  urlcolor=blue!50!black
}

\lstdefinestyle{python}{
  language=Python,
  basicstyle=\ttfamily\small,
  frame=single,
  framesep=5pt,
  breaklines=true,
  showstringspaces=false,
  columns=fullflexible
}

\title{A Note on the Characterization of EDSL in\\
\textit{Anamnesis: An Open-Source Platform for\\
Large-Scale Backstory-Conditioned Survey Simulation}}
\author{John Horton\\\small MIT Sloan \& Expected Parrot}
\date{August 25, 2026}

\begin{document}
\maketitle

\begin{abstract}
Yu et al. characterize Expected Parrot's EDSL as restricting persona
conditioning to short attribute profiles. This note clarifies that the claim
describes a default presentation rather than a limitation of EDSL. EDSL agents
can carry long narrative backstories, render them through customizable Jinja2
persona templates, and retain earlier questions and answers through survey
memory. EDSL can therefore express the backstory conditioning and sequential
context accumulation central to the Anthology procedure. Anamnesis has useful
distinctive contributions, particularly its curated persona pool,
probabilistic demographic inference and matching, and graphical workflow, but
rich narrative conditioning is not a capability absent from EDSL.
\end{abstract}

\section{The claim}

In their related-work comparison, Yu et al. write that EDSL ``provides a
Python domain-specific language for constructing AI agents with trait
dictionaries and administering surveys across multiple LLMs, but its persona
conditioning reduces to the short-attribute prompting baseline that Anthology
was designed to supersede'' \cite{yu2026anamnesis}. This characterization is
not accurate as a statement about EDSL's capabilities.

EDSL does provide a convenient dictionary of agent traits. Its default
presentation may serialize those traits compactly, which can resemble the
short demographic profiles used by some persona-prompting baselines. The
default serialization, however, is not the boundary of the abstraction.

\section{Narrative backstories in EDSL}

An EDSL agent has a customizable \texttt{traits\_presentation\_template}.
This is a Jinja2 template that can reference individual traits, the complete
trait dictionary, and an optional codebook \cite{edsl}. Consequently, a trait
can contain a multi-paragraph life history and the template can place that
history directly into the model prompt. For example:

\begin{lstlisting}[style=python]
from edsl import Agent

agent = Agent(
    traits={
        "backstory": """
        Maria grew up in ...
        [a multi-paragraph account of formative experiences,
        relationships, values, circumstances, and worldview]
        """
    },
    traits_presentation_template="""
    Adopt the perspective of the person described below.

    {{ backstory }}
    """,
)
\end{lstlisting}

There is therefore no structural distinction between an ``attribute'' holding
a short value such as an age and one holding a rich narrative. More precisely,
the presentation template is evaluated as Jinja2 and trait values are inserted
into it; the trait values need not themselves be templates. The practical
limits on backstory length are the selected model's context window, inference
cost, and experimental design---not an EDSL restriction to short profiles.

The agent's general instruction is customizable as well, providing another
way to specify how the model should inhabit or reason from the supplied life
history.

\section{Sequential survey context}

Yu et al. describe Anamnesis's Anthology answering mode as prepending a
backstory and then accumulating prior question--answer pairs so that later
answers remain consistent \cite{yu2026anamnesis}. EDSL supports this pattern
through survey memory. A researcher can expose every earlier answer to each
later question:

\begin{lstlisting}[style=python]
survey = Survey(questions).set_full_memory_mode()
results = survey.by(agent).by(model).run()
\end{lstlisting}

Alternatively, targeted memory can specify exactly which earlier answers are
included for a given question. Thus EDSL supports both full sequential context
accumulation and more controlled experimental variants.

\section{A more precise comparison}

This correction does not diminish the substantive contributions of
Anamnesis. The system described by Yu et al. combines several features into an
accessible application, including a pool of approximately 35,000 pre-generated
backstories, probabilistic demographic annotations, demographic selection and
matching procedures, and a graphical survey workflow \cite{yu2026anamnesis}.
Those are meaningful product and methodological distinctions.

The comparison would be more accurate if it separated the \emph{persona
representation} from the \emph{surrounding workflow}. For example:

\begin{quote}
EDSL defaults to dictionary-style trait presentation, but supports arbitrary
narrative backstories through long-form traits and customizable Jinja2 persona
templates, as well as sequential question--answer memory. Unlike Anamnesis,
EDSL does not provide the same pre-indexed backstory pool and probabilistic
demographic-matching workflow out of the box.
\end{quote}

The paper also suggests that EDSL requires explicit code-based specification
of scenarios, models, tools, policies, and evaluation hooks. These are optional
compositional capabilities rather than requirements for a basic survey. EDSL
also provides a command-line interface for constructing agents and surveys
without writing Python. Finally, EDSL itself is open-source software distributed
under the MIT License \cite{edsl}; it should be distinguished from any separate
hosted services offered by Expected Parrot.

\section{Conclusion}

EDSL can implement Anthology-style narrative conditioning without a change to
its agent abstraction: the backstory can be stored as a trait and rendered by
a custom persona template. Its survey-memory facilities can also reproduce
the sequential accumulation of prior responses. The statement that EDSL
``reduces to'' short-attribute prompting therefore mistakes one default
presentation for a capability limitation. The narrower and supportable contrast
is that Anamnesis packages narrative generation, demographic inference and
matching, and survey execution into a specialized graphical workflow.

\begin{thebibliography}{9}

\bibitem{yu2026anamnesis}
S.-Z. Yu, J. Suh, S. Chang, and D. M. Chan.
\newblock \emph{Anamnesis: An Open-Source Platform for Large-Scale
Backstory-Conditioned Survey Simulation}.
\newblock arXiv:2607.10628v1, 2026.
\newblock \url{https://arxiv.org/abs/2607.10628}.

\bibitem{edsl}
Expected Parrot.
\newblock \emph{EDSL: Tools for AI-powered surveys and experiments}.
\newblock Open-source software, MIT License.
\newblock \url{https://github.com/expectedparrot/edsl}.

\end{thebibliography}

\end{document}
